\section{Introduction}
\label{sec:introduction}
%==============================================================================

The development of large language models over the past years has primarily been defined by two reinforcement learning frameworks: Reinforcement Learning from Human Feedback (RLHF) \citep{ouyang2022training,christiano2017deep} and Reinforcement Learning with Verifiable Rewards (RLVR) \citep{guo2025deepseek, chen2026does}. These paradigms account for most of the improvements seen in modern instruction-following and reasoning models. Both share a basic structure where evaluative signals—whether from human raters or automated verifiers—are compressed into a scalar reward. This scalar is then used to update model weights via policy gradients. However, the information bottleneck created by reducing complex feedback to a single number is a well-documented problem that often limits the nuance a model can learn \citep{casper2023open, singhal2023long}.

In parallel, a third approach has emerged that we refer to as \textit{Verbal Reinforcement Learning} (VRL) \citep{shinn2023reflexion}. In VRL, the model is trained or steered using natural-language feedback instead of, or alongside, scalar rewards. These methods replace the compressed scalar with structured linguistic signals—such as critiques, reflections, or edit suggestions—and treat that text as the primary signal for optimization. Our position is that between 2022 and 2026, VRL shifted from a collection of isolated techniques like prompt tuning and self-correction into a core reinforcement learning paradigm that now performs on par with scalar RL in reasoning and agentic tasks.

The VRL literature currently appears fragmented because different subfields use different terminology. Researchers working on prompt optimization, generative reward models, and inference-time self-correction often work in isolation, despite using similar underlying mechanics. This survey introduces the \textit{substitution principle} to link these fields. We show that each VRL method can be defined by which part of the standard RL pipeline—the policy, reward, critic, or replay buffer—it has replaced with a language-based equivalent. When viewed this way, the commonalities are clear: the textual gradients used to optimize prompts \citep{pryzant2023protegi} are structurally very similar to the critiques used for self-correction \citep{madaan2023selfrefine} or the training signals used in semantic RL \citep{wang2025cft}.

Three main observations suggest that VRL should be treated as a single, unified field. First, unlike a scalar, a verbal critique is information-preserving; it explains \textit{why} an output is incorrect. Research shows this richer signal is more effective for tasks where the reasoning process is as important as the final result \citep{zhao2025genprm}. Second, substitution is now happening at every level of the RL stack, including the reward-design layer where models write their own reward functions in code \citep{ma2024eureka}. Third, the distinction between prompting and training is becoming less clear. Methods like feedback-conditional fine-tuning \citep{luo2025fcp} and self-rewarding language models \citep{yuan2025selfrewarding} sit on a single continuum, differing only by whether they update model weights or remain in-context.

This survey covers work where natural-language feedback is the primary optimization signal. We exclude standard scalar RLHF and basic supervised fine-tuning to focus specifically on the language feedback loop. Our contributions include a formal VRL framework that adds a language channel $\mathcal{L}$ to the standard Markov Decision Process (MDP), a taxonomy that organizes the field into four branches based on weight updates and pipeline placement, and a central thesis arguing that VRL is the primary mechanism by which modern models learn reasoning and tool-use. We suggest that future progress will come from the systematic transfer of techniques across the current divide between prompting and training.

% \paragraph{Outline.} The remainder of the paper proceeds as follows. Section~\ref{sec:formulation} formalizes VRL as an extension of the standard MDP, introduces the taxonomy (Figure~\ref{fig:taxonomy}) and the substitution principle (Table~\ref{tab:substitutions}), and justifies the structure of the survey. Section~\ref{sec:branchA} surveys non-parametric, inference-time VRL — self-refinement loops, externally grounded feedback, multi-agent collaboration, stateful memory, and tree search. Section~\ref{sec:branchC} examines verbal RL on prompts, in which the prompt itself is treated as the optimizable policy. Section~\ref{sec:branchB} turns to parametric, training-time VRL — iterative self-training, preference and alignment optimization, and feedback-conditional modeling. Section~\ref{sec:branchD} surveys verbal critics and verifiers — outcome reward models, process reward models, and generative critics — which we argue collectively function as the evaluative engine of the broader VRL ecosystem. Section~\ref{sec:discussion} synthesizes our position arguments and identifies open problems and methodological crossovers that constitute, in our view, the most promising directions for future work. Section~\ref{sec:conclusion} concludes.

% At a broad level, we taxonomize Verbal Reinforcement Learning (VRL) into two foundational domains: non-parametric and parametric, distinguished by whether the model’s underlying weights are updated during the learning process
% . Within the non-parametric umbrella, we further categorize techniques into inference-time RL and verbal RL on prompts
% . Inference-time RL focuses on episodic adaptation, where a frozen model generates an initial response, receives natural language critiques, and iteratively self-corrects its output within a single session
% . Conversely, verbal RL on prompts treats the instruction itself as an optimizable "program" or "code string," utilizing LLMs to propose and mutate linguistic instructions that trigger high performance from black-box models
% .
% On the parametric side, the field bifurcates into training-time RL and verbal critics and verifiers
% . Training-time RL focuses on skill internalization, updating model weights directly through semantic feedback to instill permanent capabilities, such as complex reasoning or alignment with human values
% . Finally, verbal critics and verifiers focus on the development of specialized "generative reward models" trained to identify and explain nuanced reasoning errors
% . These models function as the evaluative engine for VRL, addressing the challenge of scalable oversight in complex tasks where human judgment is often a bottleneck.
